\section{Theoretical Foundation}
\label{sec:theoretical_foundation}

The main idea of SBC is to use the rank statistic as a diagnostic measure.
For each simulation $i = 1, \dots, S$:

\begin{algorithm}[!ht]
    Sample $\theta_i \sim p(\theta)$\;
    Sample $y_i \sim p(y \mid \theta_i)$\;
    Obtain posterior samples $\{\theta_i^{(1)},\dots,\theta_i^{(L)}\} \gets \mathcal{A}(y_i)$\;
    Compute rank $R_i$ of $\theta_i$ among the posterior samples\;
    \caption{
        \textbf{Single iteration of the SBC procedure}.
        For each simulation, sample parameters from the prior, generate synthetic data, run inference, and compute the rank statistic.
    }
    \label{alg:sbc-iteration}
\end{algorithm}

Under the assumption that the inference algorithm, $\mathcal{A}$, correctly reproduces the posterior distribution $p(\theta \mid y)$, the rank statistic follows a discrete uniform distribution on $\{0, 1, \dots, L\}$.
Theorem \ref{thm:rank-uniformity} states and proves this result through an approach similar to the argument in Appendix B of \textcite{talts_2020}.
\begin{theorem}[Rank Uniformity Under Correct Posterior Inference]
    If the inference algorithm $\mathcal{A}$ used in Algorithm~\ref{alg:sbc-iteration} correctly reproduces the posterior distribution $p(\theta \mid y)$ for all $\theta$ and $y$, then the rank statistic $R_i$ produced by the SBC iteration is distributed uniformly on $\{0,1,\dots,L\}$.
    \label{thm:rank-uniformity}
\end{theorem}

\begin{proof}
    For simplicity, we will omit the indices in this proof.
    Now, without loss of generality, we can relabel all the posterior samples such that
    \[\theta^{(1)} \leq \theta^{(2)} \leq \dots \leq \theta^{(L)}.\]

    Then, we can write the probability mass function of the rank statistic as
    \begin{align*}
        p(r) & = \binom{L}{r} \int p(y, \theta)
                \cdot \lrp{\prod_{l=1}^{r} \int_{-\infty}^{\theta} p\lrp{\theta^{(l)}\mid \theta, y} d\theta^{(l)}}
                \cdot \lrp{\prod_{l=r+1}^{L} \int_{\theta}^{\infty} p\lrp{\theta^{(l)}\mid \theta, y} d\theta^{(l)}}
                d\theta dy \\
             & = \binom{L}{r} \int p(y, \theta)
                \cdot \lrp{\int_{-\infty}^{\theta} p\lrp{\theta^{(l)}\mid \theta, y} d\theta^{(l)}}^r
                \cdot \lrp{1 - \int_{-\infty}^{\theta} p\lrp{\theta^{(l)}\mid \theta, y} d\theta^{(l)}}^{L-r}
                d\theta dy.
    \end{align*}

    Since we condition on the data $y$, the posterior sample is independent of the generative parameter $\theta$, so
    \[p\lrp{\theta^{(l)}\mid \theta, y} = p\lrp{\theta^{(l)}\mid y}.\]

    Also, using that $p(y, \theta) = p(\theta \mid y)\cdot p(y)$ we can rewrite
    \begin{align*}
        p(r) & = \binom{L}{r} \int p(y, \theta)
                \cdot \lrp{\int_{-\infty}^{\theta} p\lrp{\theta^{(l)}\mid \theta, y} d\theta^{(l)}}^r
                \cdot \lrp{1 - \int_{-\infty}^{\theta} p\lrp{\theta^{(l)}\mid \theta, y} d\theta^{(l)}}^{L-r}
                d\theta dy \\
             & = \binom{L}{r} \int p(y) \int p(\theta \mid y)
                \cdot \lrp{\int_{-\infty}^{\theta} p\lrp{\theta^{(l)}\mid y} d\theta^{(l)}}^r
                \cdot \lrp{1 - \int_{-\infty}^{\theta} p\lrp{\theta^{(l)}\mid y} d\theta^{(l)}}^{L-r}
                d\theta dy.
    \end{align*}

    Since the algorithm correctly reproduces the posterior distribution, we have
    \[p\lrp{\theta^{(l)} \mid y} = p(\theta \mid y).\]

    Now, define
    \[u = \int_{-\infty}^{\theta} p(t \mid y) dt = F_y(\theta) \implies u \in [0, 1] \land du = p(\theta \mid y) d\theta,\]
    \noindent which allows us to change the variables in the integral. Then
    \begin{align*}
        p(r) & = \binom{L}{r} \int p(y) \cdot \int_0^1 u^r \cdot \lrp{1 - u}^{L-r} du dy \\
             & = \dfrac{L!}{r! \cdot (L - r)!} \int p(y) \cdot \dfrac{r! \cdot (L - r)!}{(L + 1)!} dy \\
             & = \dfrac{1}{L + 1} \int p(y) dy \\
             & = \dfrac{1}{L + 1},
    \end{align*}

    \noindent which is consistent with a uniform distribution over the $L + 1$ possible ranks of the set $\{0, 1, \dots, L\}$.
\end{proof}

In practice, Theorem \ref{thm:rank-uniformity} states that by repeating this process across $S$ simulations, we can plot and see visually or test statistically whether the computed ranks are uniformly distributed to validate the calibration of the inference process.
Figure \ref{fig:sbc_workflow} shows the workflow of this validation.

\begin{figure}[!ht]
    \centering
    \begin{tikzpicture}[
        node distance=0.9cm,
        >=Latex,
        thick,
        process/.style={
            draw,
            rounded corners,
            minimum width=4cm,
            minimum height=1cm,
            align=center,
            font=\small,
            fill=#1!15
        },
        process/.default=blue,
        decision/.style={
            diamond,
            aspect=2.2,
            draw,
            align=center,
            inner sep=1pt,
            fill=yellow!15,
            font=\small
        }
    ]

        % ---- Nodes ----
        \node[process=magenta]
            (start)  {Set $i  = 0$};
        \node[process=blue,   below=of start]
            (sample)  {Update i = i + 1 \\ Sample $\theta_i \sim p(\theta)$};
        \node[process=green,  below=of sample]
            (simulate) {Simulate $\boldsymbol{y_i} \sim p(\boldsymbol{y} \mid \theta_i)$};
        \node[process=orange, below=of simulate]
            (infer) {Infer $p(\theta \mid \boldsymbol{y_i})$};
        \node[process=red,    below=of infer]
            (rank) {Compute rank of $\theta_i$};
        \node[process=cyan,   below=of rank]
            (store) {Store rank $r_i$};
        \node[decision, below=of store]
            (check) {$i = S$ ?};
        \node[process=gray,   right=2cm of check, text width=4.2cm]
            (aggregate) {Check uniformity \\ of $\{r_i\}_{i=1}^S$};

        % ---- Arrows ----
        \draw[->] (start) -- (sample);
        \draw[->] (sample) -- (simulate);
        \draw[->] (simulate) -- (infer);
        \draw[->] (infer) -- (rank);
        \draw[->] (rank) -- (store);
        \draw[->] (store) -- (check);

        % Decision arrows
        \draw[->] (check.east) -- node[above]{yes} (aggregate.west);
        \draw[->, bend left=40] (check.west) to node[left]{no} (sample.west);

    \end{tikzpicture}
    \caption{
        \textbf{SBC workflow}.
        For each simulated dataset, we generate $\theta_i$, simulate data $y_i$, infer the posterior, and check whether ranks of $\theta_i$ are uniformly distributed.
    }
    \label{fig:sbc_workflow}
\end{figure}
